% Imporditud: tools/import_eksperimendid.py
% Allikas: Eesti füüsikaolümpiaadi eksperimendi ülesannete
%   kogu (Rael Kalda, 2023), ülesanne Ü73, lahendus L73.
\setAuthor{Kaur Aare Saar}
\setRound{piirkonnavoor}
\setYear{2021}
\setNumber{P E2}
\setDifficulty{2}
\setTopic{staatika}
\setVahendid{20-eurosendine münt, A4 paber pindtihedusega 80 g m$^{-2}$, joonlaud}
\setPunktid{12}
\setAllikas{Eksperimendi kogu Ü73, lahendus lk 66}
\setLahendusseis{toores}

\prob{Münt}
Leidke 20-eurosendise mündi mass.

\hint

\solu
Mõõdame joonlauaga paberi mõõtmed a = 21,0 cm ja b = 29,7 cm (1 p, kui kumbki

ei erine rohkem kui 2 mm tegelikust küljepikkusest).

Leiame paberi massi. mpaber = $\rho$ab = 5 g (1 p, kui mass ei erine rohkem kui 0,1 g. (Märkus: A4 paberi jaoks võib kasutada ka pindalavalemit A = 2$-$4m2 = 1/16 m$^2$)

Voldime paberit mööda pikemat külge mõned korrad ning teeme sellest kangi ja

kasutame mündi massi leidmiseks kangkaalu põhimõtet (1 p). Et paber ei läheks

liiga paksuks, et selle pealmine külg ei püsiks horisontaalsena, võib viimase volti-

mise jätta täisnurga alla. See suurendab ka paberi jäikust.

Paigutame mündi kangi ühte otsa (nii et pool münti jääb paberile ja pool üle paberi

ääre) ning leiame kui palju võib maksimaalselt lükata paberit üle laua serva, et see

ümber ei kukuks (2 p).

Teeme vähemalt 3 mõõtmist, ning leiame, et keskmiselt on laua peale jääva osa

pikkus x = 6,9 cm (2 p või 1 p kui on tehtud 2 või vähem mõõtmist).

Leiame paberi masskeskme kauguse laua servast. y = b $-$ x = 7,9 cm (1 p). Paneme kirja kangi reegli laua serva kui toetuspunkti suhtes (3 p õige võrrandi

eest): ympaber = xmmünt. ympaber Seega mmünt = x = 5,74 g (1 p iga vastuse eest vahemikus 5,5 g kuni 6 g).
\probend
